\documentclass[11pt]{article}
\usepackage[a4paper, margin=2.5cm]{geometry}
\usepackage[utf8]{inputenc}
\usepackage[T1]{fontenc}
\usepackage[slovene]{babel}
\usepackage{amsmath, amsthm, amssymb}
\usepackage{graphicx}

\begin{document}

\title{Brownovo gibanje}
\author{Matej Rojec}
\date {}

{\theoremstyle{definition}
\newtheorem{definicija}{Definicija}
}

{\theoremstyle{plain}
\newtheorem{izrek}{Izrek}
}
\maketitle



Brownovo gibanje (več v \([1]\) je intuitivno slučajen proces, % Sklic na knjigo
ki predstavlja naključno gibanje delcev v mediju.
    
    \includegraphics[scale=0.3]{PerrinPlot2.pdf}
    % Napis pod sliko: 
    %\caption{Reprodukcija slike iz Jean Baptiste Perrin, \emph{Mouvement brownien et réalité moléculaire}, Ann. de Chimie et de Physique (VIII) 18, 5-114, 1909}

    % Začetek definicije
    \begin{definicija}
	Standardno Brownovo gibanje \(\{B_t\}_{t \geq 0}\) je slučajen proces z naslednjimi lastnostmi: 
        \(B_0 = 0\).
        Prirastki \(B_{t_n} - B_{t_{n-1}}, B_{t_{n-1}} - B_{t_{n-2}}, \ldots, B_{t_2} - B_{t_1}, B_{t_1} - B_{t_0}\) so neodvisne slučajne spremenljivke, za vsak \(t_0 \leq t_1 \leq \cdots \leq t_{n-1} \leq t_n\).
        Za vsak \(t \geq 0\) in \(h > 0\) velja \(B_{t+h} - B_t \sim \mathcal{N}(0, h)\).
        Funkcija \(t \mapsto B_t\) je zvezna skoraj gotovo.
    \end{definicija}
    % Konec definicije
    
    Preden zapišemo izrek, definirajmo še pojem časa ustavljanja.
    
    % Začetek definicije
    \begin{definicija}
	Slučajna spremenljivka \(\tau\) na verjetnostnem prostoru \((\rho,\mathcal{F},\mathcal{P})\) z vrednostmi v %\(\)
    je čas ustavljanja glede na filtracijo !!, če velja: !!.
\end{definicija}
    % Konec definicije
    
    Zdaj lahko zapišemo izrek \(1\). % Sklic na izrek z oznako thm:stopped_brownian
    
    % Začetek izreka
    \begin{izrek}
        \label{thm:stopped_brownian}
	Naj bo \(\{B_t\}_{t \geq 0}\) (standardno) Brownovo gibanje, \(\tau\) čas ustavljanja glede na
    !! in naj velja !!.
    Potem je tudi proces:
    \[
    \hat{B} := \{B_{\tau+t} - B_\tau \mid t \geq 0\}
    \]
    (standardno) Brownovo gibanje in neodvisen od \(\mathcal{F}_\mathcal{r}\).
\end{izrek}
    % Konec izreka
    
\end{document}